\chapter{SUSTAINABILITY, INCLUSIVITY, AND SOCIETAL IMPACT}
\label{ch:sustainability}

% Short reflection (one or two paragraphs is enough per template). SDG mapping
% optional. The poster grading form awards up to +10% bonus for this, so worth
% getting right.

This project was designed to be modest in its use of resources. The quantum
engines ran on a CPU-based circuit simulator over a deliberately small dataset of
twenty images, rather than on large-scale GPU training or extended
quantum-hardware time, so the entire benchmark sweep completes on a single
workstation; the one run on real quantum hardware consumed eighty seconds of a
free-tier quota. The data raises no privacy concerns, as Flickr30k is a public
research dataset of captioned images and contains no personal identifying
information that we collect or process. The work is openly reproducible: the
source code, the benchmark harness, and the dataset configuration are publicly
available, and the harness is deterministic, so the reported results can be
regenerated from scratch.

Beyond these practical points, the study makes a small contribution to how
quantum computing is discussed. Public accounts often present quantum advantage
as imminent and general; by measuring where quantum search helps and where it
does not, and by being explicit about the data-loading cost that the headline
speed-ups omit, the work offers a counterweight to that framing. In the language
of the United Nations Sustainable Development Goals, the project speaks most
directly to SDG~4 (quality education), as an openly available teaching artefact
for quantum and classical retrieval, and to SDG~9 (industry, innovation, and
infrastructure), as an honest assessment of a candidate technology's readiness.
